
We distinguish two cases.

\medskip\noindent\textit{First case.}\quad
\(\deg Q=\deg P>p^q\).  Then \(\deg P_1<\deg P\), and the
induction hypothesis applies.

\medskip\noindent\textit{Second case.}\quad
\(\deg P=p^q\), with \(a_q\ne0\).  Then
\[
  \deg P_1=\deg P=p^q
  \qquad\text{and}\qquad
  \deg Q_1<p^q.
\]
We cannot have \(Q=0\), for otherwise \(G_1\) would not be
connected.  If the polynomial \(Q_1\) has no linear part, the
preceding transformation may be repeated.  Continuing in this
fashion, we eventually obtain an equation of the form
\[
  \tag{\(*\)\(_s\)}\label{sga3:xvii:lemma:4-1-5:equation:star-s}
  P_s(X)=-X_1+\sum_{j=1}^{q}a_j^{1/p^s}X_1^{p^j}+Q_s(X),
\]
\SGARefTarget{sga3:xvii:lemma:4:1:5:eq:star:s}
where
\[
  Q_s(X)=Q(X_2,\ldots,X_n)^{1/p^s}
\]
is an additive polynomial with a nonzero linear part, and moreover
\(\deg Q_s<p^q\).  Suppose, for example, that the coefficient of
\(X_2\) in \(Q_s\) is nonzero, and let \(-L\) be the linear part
of \(P_s(X)\).  After a linear change of coordinates, the equation
of \(G_1\) becomes
\[
  P'(X)=-L+\sum_{j=1}^{q}a_j^{1/p^s}X_1^{p^j}
    +Q'(L,X_3,\ldots,X_n),
\]
where \(Q'\) is an additive polynomial without a linear part and
\(\deg Q'<p^q\).  We are thereby reduced to the first
case.

